%%=============================================================================
%% Inleiding
%%=============================================================================

\chapter{\IfLanguageName{dutch}{Inleiding}{Introduction}}%
\label{ch:inleiding}

Voetbalclubs verzamelen vandaag meer data dan ooit tevoren. Naast klassieke wedstrijdstatistieken gaat het om GPS- en trackinggegevens, trainingsbelasting, medische of fysieke meetwaarden en gedetailleerde eventdata (passen, duels, schoten, \dots).
In theorie kan deze datastroom coaches en staf helpen om beslissingen te onderbouwen rond prestatieoptimalisatie, blessurepreventie en tactische evaluatie.
In de praktijk blijkt echter dat het louter verzamelen van data niet automatisch leidt tot betere prestaties: de meerwaarde ontstaat pas wanneer data op een gestructureerde manier worden beheerd, geanalyseerd en vertaald naar bruikbare inzichten voor de dagelijkse werking van een club~\autocite{Alamar2013}.

Tegelijk is er een duidelijk spanningsveld tussen (1) de groeiende beschikbaarheid van datatools en analysetechnieken en (2) de realiteit binnen veel (kleine tot middelgrote) voetbalomgevingen, waar tijd, middelen en gespecialiseerde data-expertise beperkt zijn.
Daardoor blijven kansen liggen: clubs beschikken vaak wel over relevante data, maar benutten die onvoldoende of inconsistent, waardoor beslissingen alsnog vooral op intuïtie, fragmentaire informatie of niet-gestandaardiseerde analyses steunen~\autocite{Rein2016}.

Deze bachelorproef vertrekt vanuit deze vaststelling en onderzoekt hoe voetbalclubs hun \emph{bestaande} data optimaler kunnen inzetten. De focus ligt dus niet op "meer data verzamelen", maar op het professioneler organiseren en toepassen van data die al beschikbaar is (of realistisch beschikbaar kan zijn).
Het einddoel is een praktisch bruikbaar stappenplan en een proof-of-concept (PoC) dat toont hoe inzichten op een reproduceerbare en toegankelijke manier gegenereerd kunnen worden.

% \begin{itemize}
%   \item context, achtergrond
%   \item afbakenen van het onderwerp
%   \item verantwoording van het onderwerp, methodologie
%   \item probleemstelling
%   \item onderzoeksdoelstelling
%   \item onderzoeksvraag
%   \item \ldots
% \end{itemize}

\section{\IfLanguageName{dutch}{Probleemstelling}{Problem Statement}}%
\label{sec:probleemstelling}

Veel voetbalclubs hebben de voorbije jaren geïnvesteerd in databronnen (bijvoorbeeld wedstrijd- en trainingsstatistieken, GPS, tracking, medische opvolging), maar halen onvoldoende rendement uit die informatie.
Deze onderbenutting heeft zowel sportieve als organisatorische gevolgen: trainings- en wedstrijdkeuzes worden minder onderbouwd, signalen rond belasting of prestatie-evolutie worden mogelijk te laat opgemerkt, en data blijft verspreid over verschillende tools, bestanden of personen, waardoor kennis niet schaalbaar of overdraagbaar is~\autocite{Rein2016}.

De doelgroep van deze bachelorproef bestaat uit clubs die reeds data verzamelen of data ter beschikking hebben, maar moeilijkheden ervaren bij het verwerken, analyseren en interpreteren ervan.
Het gaat in het bijzonder om clubs zonder uitgebouwde data-afdeling of zonder fulltime data-analist, waar de implementatiedrempel (tijd, expertise, tooling) hoog blijft~\autocite{Alamar2013}.

\section{\IfLanguageName{dutch}{Onderzoeksvraag}{Research question}}%
\label{sec:onderzoeksvraag}

De centrale onderzoeksvraag luidt:

\begin{quote}
  \textbf{Hoe kunnen voetbalclubs hun bestaande data optimaler inzetten om de kwaliteit van besluitvorming en sportprestaties te verbeteren?}
\end{quote}

Om deze onderzoeksvraag behapbaar te maken, worden de volgende deelvragen onderzocht:

\begin{enumerate}
  \item \textbf{Welke soorten data verzamelen voetbalclubs vandaag en hoe worden die momenteel gebruikt?}
  \item \textbf{Welke knelpunten ervaren clubs bij het verwerken, analyseren en interpreteren van die data (technisch en organisatorisch)?}
  \item \textbf{Welke analysetechnieken en tools zijn haalbaar en bruikbaar voor clubs met beperkte data-expertise?}
  \item \textbf{Welke succesfactoren bepalen of datagedreven werken duurzaam ingebed raakt in processen en rollen binnen een club?}
\end{enumerate}

\section{\IfLanguageName{dutch}{Onderzoeksdoelstelling}{Research objective}}%
\label{sec:onderzoeksdoelstelling}

De doelstelling van deze bachelorproef is tweeledig:

\begin{itemize}
  \item \textbf{Praktisch:} een concreet en toepasbaar stappenplan uitwerken waarmee voetbalclubs hun datagebruik gefaseerd kunnen professionaliseren (databeheer, toolkeuzes, rollen, implementatievolgorde).
  \item \textbf{Technisch:} een \textbf{proof-of-concept} ontwikkelen (in Python) die aantoont hoe een selectie van relevante voetbaldata kan worden opgeschoond, gestructureerd, geanalyseerd en gevisualiseerd met interpreteerbare technieken (bv.\ beschrijvende analyses, eenvoudige modellen, clustering of anomaliedetectie), zodat inzichten reproduceerbaar worden.
\end{itemize}

Succes wordt beoordeeld op basis van (1) de mate waarin het stappenplan concreet, gefaseerd en realistisch toepasbaar is in clubs met beperkte middelen, en (2) de mate waarin de PoC transparant, reproduceerbaar en begrijpelijk is voor niet-specialisten (met een duidelijke output, visualisaties en documentatie).

\section{\IfLanguageName{dutch}{Opzet van deze bachelorproef}{Structure of this bachelor thesis}}%
\label{sec:opzet-bachelorproef}

% Het is gebruikelijk aan het einde van de inleiding een overzicht te
% geven van de opbouw van de rest van de tekst. Deze sectie bevat al een aanzet
% die je kan aanvullen/aanpassen in functie van je eigen tekst.

De rest van deze bachelorproef is als volgt opgebouwd:

In Hoofdstuk~\ref{ch:stand-van-zaken} wordt de stand van zaken besproken rond voetbaldata, sportanalytics en datagedreven besluitvorming, met aandacht voor datatypes, analysetechnieken, tooling en implementatie-uitdagingen.

In Hoofdstuk~\ref{ch:methodologie} wordt de onderzoeksopzet toegelicht: hoe de literatuurstudie, praktijkbevraging (interviews/gesprekken), requirements- en knelpuntanalyse, de ontwikkeling van de proof-of-concept en de synthese naar een stappenplan worden uitgevoerd.

% TODO: Vul hier aan voor je eigen hoofstukken, één of twee zinnen per hoofdstuk

In Hoofdstuk~\ref{ch:conclusie} ten slotte worden de conclusies geformuleerd en wordt een antwoord gegeven op de onderzoeksvragen. Daarbij wordt ook besproken welke beperkingen er zijn en welke pistes relevant zijn voor toekomstig onderzoek.